\section{Prohibited paths and centered traces}\label{q:paths-trace}

We now construct the sparse sets used in \Cref{q:cuts}.  Their
removal will impose enough structure on closed walks to control a
matrix moment.  Throughout this section the bin index $j$ is fixed,
and
\[
 s=\lfloor L^{1/10}\rfloor,\qquad k=\lfloor L\rfloor.
\]
All assertions are for sufficiently large $L$, with the fixed integers
$h,l$ and the absolute parameters $A,W$ held fixed.

The centered graph and high-trace strategy was developed by Helfgott
and Radziwi\l\l\ \cite[Sections~2--7]{HR21}.
Pilatte developed its composite-label and nonbacktracking form
\cite[Sections~2--9]{Pilatte26}.  Here each step carries a prime tuple
and a padding divisor; the trace expansion keeps the padding weights
and the common residue coordinates explicit.

\subsection{Positive paths and prohibited words}

A signed step is a displacement
$D_i=\varepsilon_i hq_i d_i$, where $\varepsilon_i\in\{-1,1\}$
and $(d_i,q_i)\in\mathcal S_j$.
A word of $m$ steps, starting at $x$, visits the sites
$x+\sum_{a<i}D_a$ for $1\le i\le m+1$.
It is \emph{positive} if $q_id_i$ divides its departure site for
every $i$.  This condition also holds at its arrival site, so
positivity survives restriction and reversal.  These definitions
apply in the auxiliary residue model as well: divisibility then means
the prescribed equality in the corresponding residue coordinate.
A path need not remain inside any finite interval.

The next definition separates an equality pattern from an arithmetic
condition on its labels.  Within the pattern, a tuple prime may persist for
several consecutive steps, but cannot return after disappearing.
The additional condition is a divisibility relation on a suffix.

\begin{definition}\label{q:prohibited}
A numerical word is \emph{forward prohibited} if it has length
$3\le m\le s$ and satisfies
three conditions:
\begin{enumerate}
\item Consecutive whole tuples $d_i,d_{i+1}$ are unequal.
\item For every prime in $P$, its indices of appearance among
$d_1,\ldots,d_m$ form an interval, if nonempty.
\item Some prime $p\mid d_1$, absent from $d_m$, satisfies
\begin{equation}\label{q:prohibited-relation}
 p\mid\sum_{i=a}^m D_i\qquad\text{for some }1<a<m.
\end{equation}
\end{enumerate}
A forward prohibited word is \emph{minimal} if no shorter contiguous
subword is forward prohibited in either orientation.  A \emph{witness}
is a positive path realizing a minimal forward prohibited word.
Let $\mathcal B_j$ be the set of sites at which a witness
starts, and set
$\sigma_j(n)=\1_{n\notin\mathcal B_j}$.
\end{definition}

Thus we specialize the event and indicator used in
\Cref{q:tau,q:cuts}.  Minimality is a condition on the numerical word,
whereas positivity is a condition on the residue coordinates at its
starting site.  Any positive prohibited word
contains a witness starting at one of its vertices: repeatedly choose
a shorter prohibited subword, reversing when necessary.  Length
strictly decreases, while positivity is preserved.
The structural use of this deletion appears in
\Cref{q:block-intervals}: on surviving positive paths of length at most
$s$, with unequal consecutive tuples, every tuple prime has consecutive
occurrences and no nonempty subinterval has zero displacement.

We first give two elementary counting facts, including the order of
summation needed later.  A prime slot is an occurrence of a prime in
one of the tuples or paddings.  An equality pattern partitions these
slots into classes, each class representing one prime variable; tuple
classes also record their column $P_i$.

\begin{lemma}[Reciprocal counting and elimination]\label{q:crude-count}
Consider at most $R\le4L$ signed steps together with additional
records having at most $L^{O(1)}$ choices per prime slot.  Then the number of
equality patterns and such records is at most
$\exp(O(R\log^2L))$.  With one reciprocal for every distinct prime,
the sum over their numerical values has the same bound, also after
multiplication by $Lu(q_i)$ for each step.

Suppose, with the padding values fixed, that a recorded relation in
the tuple primes is
\[
 az+c\equiv0\pmod p,\qquad a\not\equiv0\pmod p,
\]
where $a,c,p$ are independent of the selected prime variable $z$.
Its reciprocal sum costs at most $(1+L)/H_0$ in place of an
unrestricted prime sum.  Such savings multiply when selected variables
are ordered so that each relation uses only its own variable, earlier
selected variables, and nonselected variables.
\end{lemma}

\begin{proof}
There are $O(R\log L)$ slots, since each tuple has $J=O(\log L)$
factors and every eligible padding has at most $100\log L$ factors.
A partition of $N$ slots has at most $N^N$ descriptions.  Signs,
factor counts and all recorded indices therefore cost
$\exp(O(R\log^2L))$.  Each numerical class has reciprocal mass
$O(\log L)$, by the prime harmonic bounds, and
$Lu(q_i)\le L^{1+100\log4}$.  This proves the first assertion.
Dropping bin or numerical distinctness restrictions enlarges these
positive sums.  If distinct abstract classes are assigned the same
numerical prime in such an enlargement, they retain their separate
reciprocal factors.

For the second assertion the coefficient is invertible modulo $p$,
so $z$ lies in one residue class.  Enlarge from primes to integers and
compare the reciprocal sum along that progression with its integral:
\[
 \sum_{\substack{H_0\le n\le e^L\\n\equiv c'\!\!\pmod p}}\frac1n
 \le\frac1{H_0}+\frac Lp\le\frac{1+L}{H_0}.
\]
Fix the nonselected variables and sum the selected ones in reverse
order.  Every remaining earlier relation is independent of the
variable currently being summed.  Its own coefficient test is retained;
if this test fails, the admissible inner sum is zero.  The displayed
bound is uniform in the remaining variables, and backward induction
proves multiplication of the savings.  After these sums, the other
prime variables are bounded by the first assertion.
\end{proof}

\begin{lemma}[Density and complexity of the deleted event]\label{q:bad}
The event in \Cref{q:prohibited} is translation covariant and satisfies
\[
 \Pbb(n\in\mathcal B_j)
 \le\exp\bigl(-\tfrac12L^{1-\delta}\bigr).
\]
As a function of the residue equalities it has a Boolean circuit of
depth two and size at most $\exp(L^3)$ for sufficiently large $L$.
\end{lemma}

\begin{proof}
We may count all forward prohibited positive words, without requiring
minimality.  For a fixed numerical word, positivity either gives
inconsistent residue conditions or fixes one residue at each distinct
prime in its steps.  Its probability is accordingly zero or the product
of their reciprocals.

At the last transition choose a column whose prime changes.  Its last
prime $z$ occurs nowhere earlier, by the interval condition.  In
\eqref{q:prohibited-relation}, the contribution involving $z$ is
exactly the last step.  Thus its coefficient is
$\varepsilon_m hq_m d_m/z$.  The controlling prime $p$ is absent from
$d_m$, from $h$ and from every padding, so this coefficient is a unit
modulo $p$.  It is independent of $z$.  Fixing the equality pattern,
padding and other tuple values, \Cref{q:crude-count} supplies one
saving $(1+L)/H_0$.  The remaining count, summed over $3\le m\le s$,
is at most
\[
 \frac{1+L}{H_0}\exp(O(s\log^2L))
 \le\exp\bigl(-\tfrac12L^{1-\delta}\bigr).
\]
This proves the probability bound.

Each eligible numerical pair satisfies $dq\le\exp(100L+1)$.
The disjoint prime supports recover $d$ and $q$ uniquely from their
product.  Thus the number of numerical words of length at most $s$
is $\exp(O(sL))$.  Preselect those satisfying the numerical prohibition
and minimality conditions.  Positivity of one such word is a conjunction
of $O(s\log L)$ residue equalities, one per prime occurrence.  Taking their
disjunction gives the stated depth and size bounds.  Translation merely
shifts the same residue coordinates.  This verifies all the hypotheses
on the deleted events in \Cref{q:cuts}.
\end{proof}

\subsection{The matrices and their closed words}

We use the cutoff $\tau_j(n,d)$ of \Cref{q:tau}: it includes the
$Q$-degree cutoff, the bound on $\rho_j$, and $\sigma_j$, but omits
the $P$-degree cutoff.  Let
\[
 M=\lceil e^{103L}\rceil,
\]
and consider the sites $t+1,\ldots,t+M$.  For each $d\in\mathcal D$
define a real symmetric matrix $A_d$.  For an increasing edge
$n'=n+hqd$ within the block, $(d,q)\in\mathcal S_j$, its entry is
\begin{equation}\label{q:matrix-entry}
 A_d(n,n')=
 \frac{Lu(q)\tau_j(n,d)\tau_j(n',d)}{\sqrt{w(n)w(n')}}
 \1_{q\mid n}\1_{n\equiv bqd\!\!\pmod l}
 \prod_{p\mid d}\left(\1_{p\mid n}-\frac1p\right).
\end{equation}
Set $A_d(n',n)=A_d(n,n')$, and set all other entries to zero.
The increasing edge determines $q$ uniquely.  Let $E$ be the
orthogonal coordinate projection to sites satisfying
$\omega_P(n)\le6WJ$.
Multiplying a row entry by $\sqrt{w(n')}/\sqrt{w(n)}$ changes its
padding weight to $Lu(q)/w(n)$; in each fixed orientation, the sum of
these padding weights, with their divisibility and bin conditions,
is at most $L\rho_j(n,d)\le K$ at a retained site.
On positive integer blocks, testing an increasing edge against $\sqrt w\,\lambda$
cancels both square-root weights, and imposing $E$ at both endpoints
recovers $L$ times the retained centered edge weight in
\eqref{q:cut-averages}.
On the direct sum of the $d$-indexed copies
of the block space define
\[
 H_{d,d'}=\1_{d\ne d'}A_{d'}.
\]
This matrix need not be self-adjoint.  Its powers record paths with
unequal consecutive whole tuples, even if their paddings differ.

\begin{theorem}[The nonbacktracking moment]\label{q:trace}
There is an absolute constant $C_2$ such that
\begin{equation}\label{q:trace-bound}
 \E\|H^k\|_{\mathrm{HS}}^2
 \le\bigl[K(C_2\sqrt W)^J\bigr]^{2k}.
\end{equation}
The norm is the unnormalized Hilbert--Schmidt norm.  The same estimate,
after increasing $C_2$ absolutely, holds when block origins are averaged
over any integer interval of length at least $\exp(\tfrac12L^A)$.
The threshold for $L$ may depend on the fixed data.
\end{theorem}

We reduce this statement to positive counts of closed words in the
rest of this section.  The next two sections dispose of words with
many arithmetic restrictions and count the remaining patterns.

An entry of $H^k$ expands over successive copy labels
$d_0,\ldots,d_k$ with $d_{i-1}\ne d_i$ and the corresponding products
of matrix entries.  Squaring and summing joins two site paths with
common endpoints.  Reverse the second using symmetry of $A_d$.
The result is a closed word of $2k$ steps, with unequal consecutive
tuples within each half.  No such condition is needed across either
join.  Taking absolute values after each word's expectation bounds
the moment by the resulting sum.  The external copy indices and initial
site cost at most
\begin{equation}\label{q:external-dimension}
 |\mathcal D|^2M\le\exp(107L+O(1)).
\end{equation}
Indeed $d\le e^{2L}$.  Once the word is recorded, there is no copy-index
choice at each step beyond its recorded tuple.

The finite-origin comparison is applied to this signed expansion,
before taking absolute values of word expectations. Fix a numerical
word, its initial block index, and its external copy indices. All
visited sites are then fixed translates of the block origin. At each
site only $Q$-divisor sets of size at most $m_Q=\lfloor400\log L\rfloor$
survive. By \eqref{q:q-state-count}, enumerating the exact sets at all
$2k$ departures costs at most $\exp(O(kL\log L))$ states. These events
specify absence of all other $Q$ primes, so the weight denominators and
the $Q$-dependent cutoffs become scalars. Some state conjunctions may
be inconsistent; their probability is then zero. No independence
between different sites is used.

Expand the at most $2kJ$ centered factors into indicators and scalars,
retaining their signs. This gives at most $2^{2kJ}$ terms with
coefficient magnitudes at most one. Only prime occurrences in the main
word enter this expansion; no complete $P$-divisibility state is listed.
Every other scalar factor per step has magnitude at most $Lu(q)$,
since $w\ge1$ and the cutoffs are at most one. As
$u(q)\le L^{100\log4}$, the product of these factors is
$\exp(O(k\log L))$.

The remaining events are prime divisibility tests, progression tests,
the exact $Q$ states, and absence of the prohibited paths from
\Cref{q:bad}. The latter events have depth two before complementation
and size at most $\exp(L^3)$. Conjoining the tests at all $2k$ sites
adds a single AND level, so their combined depth is at most $20$ and
size at most $\exp(L^6)$ for large $L$. The $Q$-degree and
$\rho_j$ cutoffs are already decisions on the exact states. Block
restrictions are numerical conditions and add no circuit levels.

Each eligible pair has $dq<e^{100L+\eta}\le e^{101L}$, and disjointness
of $P$ and $Q$ recovers $(d,q)$ from this product. The number of signed
numerical words of length at most $2k$ is therefore at most
$(2k+1)(2e^{101L})^{2k}$. The initial block index and two external
copy indices contribute the factor $M|\mathcal D|^2=\exp(O(L))$.
Combining these counts, the total absolute coefficient sum in all the
indicator expansions is at most
\begin{equation}\label{q:word-comparison-budget}
 \exp(O(kL\log L))=\exp(O(L^2\log L))
 \le\exp(O(L^5)).
\end{equation}
Thus \Cref{q:finite-expansion} compares the entire signed sum with
total error at most $\exp(-L^9)$. We may then take the absolute values
of its product-law word expectations. It suffices to prove the
auxiliary-product estimate; the integer estimate follows with this
negligible error. The comparison uses the truncated $Q$ states,
without requiring an integer mean for the full weight $w^2$.

\subsection{Integrating the centered prime coordinates}

Fix a numerical closed word and write
\[
 x+\xi_i\quad\text{for its departures},\qquad
 \xi_i=\sum_{a<i}D_a\quad(1\le i\le2k).
\]
Its contribution has the form
\begin{equation}\label{q:word-product}
 R\,G\prod_{i=1}^{2k}\prod_{p\mid d_i}
       \left(\1_{p\mid x+\xi_i}-\frac1p\right),
 \qquad G=\prod_{i=1}^{2k}\sigma_j(x+\xi_i).
\end{equation}
Here $R\ge0$ is independent of all $P$-residue draws.  Its weight part
is exactly
\begin{equation}\label{q:departure-weights}
 \prod_{i=1}^{2k}
 \frac{Lu(q_i)}{w(x+\xi_i)}\1_{q_i\mid x+\xi_i},
\end{equation}
multiplied by the $Q$-degree and $\rho_j$ cutoffs at incident endpoints,
the orientation-dependent progression tests, and block restrictions.
The denominator identity follows from closure, counting every visit
with multiplicity.  Each centered factor agrees at the two endpoints
of its edge.  For a reversed edge its progression test is still read
at the lower endpoint, and remains in $R$.
Although $R$ can depend on the numerical $P$ labels, it does not depend
on their residue draws.  This is why the $P$-degree projection has not
yet been imposed on the matrices.

Call a $P$ prime appearing at exactly one step of the word a
\emph{singleton}.  Let $\mathcal U$ be their set and $S_1=|\mathcal U|$.
If $G$ were independent of one singleton coordinate, centering in that
coordinate would make the word expectation zero.  We retain this
cancellation while accounting for the dependence introduced by the
vertex deletions.

At every occurrence
of a nonsingleton expand the centered factor into its indicator and
its scalar $-1/p$.  Call these choices \emph{lit} and \emph{unlit},
respectively, and let $U$ be the number of unlit occurrences.
An unlit choice imposes no nondivisibility condition.
For a prime $p$, write $\ell_p,u_p$ for its lit and unlit counts.
Lit consistency means that all its lit departure offsets are congruent
modulo $p$.  Define
\[
 b_{\mathcal L}=
 \prod_{p\ \mathrm{singleton}}\frac1p
 \prod_{p\ \mathrm{nonsingleton}}p^{-u_p-\1_{\ell_p>0}}.
\]

If lit consistency fails, the designated term vanishes; set
$\Delta G=0$ in this case.  Otherwise, starting from one common draw
of the prime residues, overwrite every
nonsingleton with a lit occurrence by its forced residue.  For each
$H\subseteq\mathcal U$, also overwrite the singleton coordinates in
$H$ by the residues forced at their sole occurrences, leaving the
other singleton coordinates at their original values.  Call the
resulting configuration a \emph{hybrid}, and let $G_H$ be the value of
$G$ there.  The mixed difference, forced minus original in each
singleton coordinate, is
\[
 \Delta G=\sum_{H\subseteq\mathcal U}(-1)^{S_1-|H|}G_H.
\]

\begin{lemma}[Exact centering before absolute values]\label{q:mixed-difference}
The absolute expectation of a designated term of
\eqref{q:word-product} is at most
\begin{equation}\label{q:mixed-majorant}
 \1_{\mathrm{lit\ consistency}}\,b_{\mathcal L}
       \E\bigl[R|\Delta G|\bigr].
\end{equation}
In particular $|\Delta G|\le2^{S_1}$.
\end{lemma}

\begin{proof}
For a uniform residue $Z$ modulo $p$, a fixed residue $a$, and any
function $F$ of that coordinate,
\[
 \E\left[\left(\1_{Z=a}-\frac1p\right)F(Z)\right]
 =\frac1p\E\bigl[F(a)-F(Z)\bigr].
\]
Apply this identity successively to the independent singleton
coordinates.  At a nonsingleton, incompatible lit tests give zero;
otherwise they force one residue and supply $1/p$, while every unlit
choice supplies the scalar $-1/p$.  These operations produce exactly
the factor $b_{\mathcal L}$ and the mixed difference, up to the signs
of the unlit scalars.  The factor $R$ is unchanged in every operation.
Taking the absolute value after these integrations proves the bound.
The final assertion follows because $G$ takes values in $\{0,1\}$.
\end{proof}

The square-root dependence on $W$ in \eqref{q:trace-bound} comes
from repeated prime labels.  Among $2kJ$ tuple-prime occurrences there
are at most $Jk+S_1/2$ distinct labels, since every nonsingleton occurs
at least twice.  For a fixed equality pattern, summing one reciprocal
per label therefore costs at most $(2W)^{Jk+S_1/2}$.
To use this count, we must still control the number of patterns and
sum the padding weights uniformly.  We first discard terms where
extra unlit reciprocals, independent lit-consistency congruences, or
relations forced by many singletons give a stronger saving directly.
The forest count and uniform padding sum in \Cref{q:forest-section}
will handle the remaining terms.

For negligible classes of words we will discard denominators and
cutoffs in $R$, but retain all padding divisibility tests.
Independence then gives one reciprocal per distinct $Q$ prime.
The factor $b_{\mathcal L}$ gives at least one reciprocal per distinct
$P$ prime.  Together with \Cref{q:crude-count}, the total cost before
any additional saving is $\exp(O(L\log^2L))$, including
\eqref{q:external-dimension}, lit designations and $2^{S_1}$.
Padding values may be fixed first during a constrained sum over tuple
primes; their original bound $q_i\le e^{100L+1}$ is retained when
needed for coefficient-size estimates.

\begin{lemma}[Many unlit occurrences]\label{q:many-unlit}
The total contribution to the moment majorant from designated terms
with $U>L^{1/50}$ is at most $\exp(-L^{1+\delta})$.
\end{lemma}

\begin{proof}
If a nonsingleton has a lit occurrence, each of its unlit occurrences
supplies an extra reciprocal beyond its basic $1/p$.
If all its $m\ge2$ occurrences are unlit, it supplies $p^{-m}$ and
hence at least $m/2$ extra reciprocal powers beyond the basic one.
Thus at least $U/2$ extra powers remain, all at primes at least $H_0$.
The total is at most
\[
 \exp(O(L\log^2L))H_0^{-U/2}
 \le\exp\bigl(O(L\log^2L)-\tfrac12L^{1.015}\bigr)
 \le\exp(-L^{1+\delta}).
\]
\end{proof}

We now turn to the two remaining sources of arithmetic savings:
independent congruences and singleton primes.
